\documentclass{beamer}
\usetheme{Madrid}
\usepackage[utf8]{inputenc}
\usepackage[ukrainian]{babel}
\usepackage{amsmath}
\usepackage{hyperref}
\usepackage{algorithm}
\usepackage{algpseudocode}

\setbeamertemplate{footline}{
  \hfill%
  \usebeamercolor[fg]{page number in head/foot}%
  \usebeamerfont{page number in head/foot}%
  \insertframenumber{} / \inserttotalframenumber\hspace{1em}
}


\title[Analysis of Efficiency and Security of Existing \\ BFT Paxos-Based Algorithms]

\date{}

\author[FFFF\quad|\quad Фесенко А. В.]{
    Illia Melnyk, Oleksandr Kurbatov, Oleg Fomenko, Volodymyr Dubinin and Yaroslav Panasenko
}
\institute{
    Науковий керівник: Фесенко Андрій Вʼячеславович
}



\begin{document}

\begin{frame}
  \titlepage
\end{frame}

\begin{frame}{Вступ і актуальнiсть}

\vspace{0.5em}
\textbf{Квантові обчислення} радикально змінюють уявлення про безпеку:
\begin{itemize}
  \item Класичні алгоритми $\Rightarrow$ вразливі до квантових атак 
  \item Потрібен \textbf{перехiд} до постквантових протоколiв
\end{itemize}

\vspace{0.5em}
\textbf{Гiбриднi криптографiчнi схеми} — перехiдний варiант:
\begin{itemize}
  \item Об’єднують: класичну + PQ криптографію
  \item Зберігають зворотну сумiснiсть і надiйнiсть
\end{itemize}


\vspace{0.5em}
\begin{block}{Актуальне питання}
  \centering
  \textbf{Чи достатньо безпечним є механiзм об’єднання ключiв (CtKDF)?}
\end{block}

\end{frame}


\begin{frame}{Мета дослiдження}

\begin{block}{Мета: }
\textbf{довести, що \textbf{CtKDF} \textbf{не є IND-CPA стiйкою}, якщо функцiя створення ключа (KDF) є нестiйкою}
\end{block}


\vspace{1em}
\textbf{Завдання дослiдження:}
\begin{enumerate}
  \item Побудувати \textit{формальну iгрову модель} для атаки типу IND-CPA.
  \item Визначити \textit{перевагу зловмисника} при використаннi слабкої функції створення ключа.
  \item Довести \textit{теорему нестiйкостi} гiбридної схеми CtKDF.
\end{enumerate}

\end{frame}


\begin{frame}[fragile]{CtKDF: суть механізму}

\textbf{CtKDF} --- гібридний механізм об’єднання ключів за допомогою конкатенації:

\vspace{0.5em}
\begin{algorithm}[H]
 \caption{CtKDF(context, label, l, psk, k, P, R)} 
 \begin{algorithmic}
  \State $secret \gets psk \| k_1 \| k_2 \| \dots \| k_n$
  \State $context' \gets f(context, P, R)$
  \State \textbf{return} $KDF(secret, label, context', l)$
 \end{algorithmic}
\end{algorithm}

\vspace{0.5em}
\begin{itemize}
  \item $k_i$ --- секрети з різних схем
  \item $psk$ --- попередньо відомий ключ (pre-shared key)
  \item $f(\cdot)$ --- функція форматування контексту
  \item $KDF$ --- функція створення ключа 
\end{itemize}

\end{frame}



\begin{frame}{Уразливість CtKDF при слабкому KDF}

\begin{block}{Ключове спостереження}
  Навіть при стійких компонентах, слабка функція KDF \textbf{руйнує всю конструкцію}.
\end{block}

\vspace{0.5em}

\textbf{Чому це відбувається?}
\begin{itemize}
  \item CtKDF не маскує структуру секрету, а лише \textit{перетворює} його через KDF
  \item При слабкій функції створення ключа: вихід \textit{містить детерміновані ознаки вхідних даних}
  \item Зловмисник може \textit{відновити частини $psk$, $k_i$, $len$ тощо}
  \item Це відкриває шлях до \textit{атаки типу IND-CPA} через ігрову модель
\end{itemize}

\end{frame}


\begin{frame}[fragile]{Ігрова модель: G\textsubscript{0} проти G\textsubscript{1}}

\vspace{0.5em}
\textbf{Мета:} змоделювати здатність зловмисника розрізнити \textit{реальний} та \textit{випадковий} вихід функції створення ключа

\vspace{0.2em}

\begin{columns}
  \begin{column}{0.48\textwidth}
    \scriptsize
    \begin{block}{Гра $G_0^{\text{kdf-weak}}$}
      \vbox to 3cm {
        \begin{algorithmic}
          \State $(inp, a) \gets \mathit{Input()}$
          \If{$\perp(inp, a)$}
            \State \Return $0$
          \EndIf
          \State $s \gets \mathit{Salt()}$
          \State $(c, len) \gets \mathcal{A}(a, s)$
          \State $out \gets (inp, s, c, len)$
          \State \Return $\mathcal{A}(out)$
        \end{algorithmic}
        \vfill
      }
    \end{block}
  \end{column}

  \begin{column}{0.48\textwidth}
    \scriptsize
    \begin{block}{Гра $G_1^{\text{kdf-weak}}$}
      \vbox to 3cm {
        \begin{algorithmic}
          \State $(inp, a) \gets \mathit{Input()}$
          \If{$\perp(inp, a)$}
            \State \Return $0$
          \EndIf
          \State $s \gets \mathit{Salt()}$
          \State $(c, len) \gets \mathcal{A}(a, s)$
          \State $out \gets \{0,1\}^{len}$ \Comment{випадковий рядок}
          \State \Return $\mathcal{A}(out)$
        \end{algorithmic}
        \vfill
      }
    \end{block}
  \end{column}
\end{columns}


\vspace{0.5em}
\textbf{Ціль зловмисника:} вгадати, з якої гри походить $out$

\end{frame}



\begin{frame}{Перевага зловмисника}

\vspace{0.5em}
\textbf{Оцінка переваги зловмисника:}

\[
\text{Adv}_{\text{kdf-weak}}^{\text{KDF, Input, Salt}}(\mathcal{A}) =
\left| \Pr[G_0^{\text{kdf-weak}}(\mathcal{A}) = 1] - \Pr[G_1^{\text{kdf-weak}}(\mathcal{A}) = 1] \right|
\]

\vspace{1em}
\begin{itemize}
  \item $G_0$: зловмисник отримує реальний вихід функції створення ключа
  \item $G_1$: зловмисник отримує випадковий бітовий рядок довжини $len$
  \item $\mathcal{A}$ — стратегія зловмисника, яка повертає біт $b'$
\end{itemize}

\vspace{0.5em}
\begin{block}{Інтерпретація}
  Якщо \textbf{Adv близьке до 1}, то зловмисник легко відрізняє вихід KDF від випадкового значення.\\
  Це означає: \alert{функція є нестійкою, а схема CtKDF — уразливою}.
\end{block}

\end{frame}


\begin{frame}{Стратегія зловмисника}

\vspace{0.5em}
\textbf{Ідея атаки:} проаналізувати вихід $out$ і визначити, чи відповідає він очікуваній структурі

\vspace{0.5em}
\begin{block}{Очікувана структура виходу}
  \[
    out \approx (inp', s', c', len')
  \]
\end{block}

\vspace{0.5em}
\begin{itemize}
  \item Якщо вихід \textbf{містить зрозумілий формат} — зловмисник робить висновок: $G_0$
  \item Якщо \textbf{не має структури / виглядає випадково} — припускає: $G_1$
  \item Стратегія не потребує дешифрування або ключів, лише \textit{аналізу форми}
\end{itemize}

\vspace{0.5em}
\textbf{Результат:} \textit{перевага зловмисника наближається до 1}, якщо KDF слабка

\end{frame}


\begin{frame}{Модифікація CtKDF без KDF}

\vspace{0.5em}
\textbf{Упрощений варіант CtKDF:} KDF повністю виключено зі схеми

\vspace{0.5em}
\begin{block}{CtKDF(context, label, l, psk, k, P, R)}
\begin{algorithmic}
  \State $secret \gets psk \| k_1 \| k_2 \| \dots \| k_n$
  \State $context' \gets f(context, P, R)$
  \State \textbf{return} $secret$
\end{algorithmic}
\end{block}

\vspace{0.5em}
\begin{itemize}
  \item Жодної хеш-функції, стискання або обфускації — \textbf{вся структура зберігається}
  \item Навіть без KDF, зловмисник може \textit{відновити або розпізнати} частини $psk$, $k_i$, $len$
  \item CtKDF у такому вигляді є \textit{демонстративно нестійкою}
\end{itemize}

\vspace{0.5em}
\textbf{Призначення:} використовується як допоміжна модель у побудові атак на повну схему

\end{frame}



\begin{frame}{Гра IND-CPA для CtKDF}

\vspace{0.5em}
\textbf{Мета:} довести, що зловмисник може виграти IND-CPA гру, якщо функція створення ключа є нестійкою

\vspace{1em}
\begin{block}{Суть гри IND-CPA}
\begin{itemize}
  \item Зловмисник бачить публічні ключі: $pk_1, pk_2, \dots, pk_n$
  \item Отримує доступ до оракула інкапсуляції
  \item Надсилає два повідомлення $m_0$, $m_1$
  \item Отримує шифротекст для $m_b$, має вгадати біт $b$
\end{itemize}
\end{block}

\vspace{0.5em}
\textbf{При слабкому KDF:}
\begin{itemize}
  \item Вихід секрету зберігає структуру
  \item Можна застосувати попередню \textit{структурну стратегію атаки}
\end{itemize}

\vspace{0.5em}
Гра IND-CPA \textbf{зводиться до гри $G_0/G_1$}, отже:
\[
  Adv^{\text{ind-cpa}}_{\text{CtKDF}}(\mathcal{A}) \approx Adv^{\text{kdf-weak}}(\mathcal{A})
\]

\end{frame}


\begin{frame}{Формалізація втрати IND-CPA стійкості при слабкій фцнкції створення ключа}

\vspace{0.5em}

\vspace{1em}
\begin{block}{Теорема}
Якщо функція створення ключа $KDF$ є нестійкою, то:
\[
  Adv^{\text{ind-cpa}}_{\text{CtKDF}}(\mathcal{A}) \geq 1 - (\varepsilon_1 + \varepsilon_2)
\]
де $\varepsilon_1$, $\varepsilon_2$ — похибки при моделюванні / редукції.
\end{block}

\vspace{1em}
\begin{itemize}
  \item Навіть при IND-CPA-стійких схемах, схема CtKDF може бути \textit{вразливою загалом}
  \item CtKDF не успадковує безпеку своїх компонент без належної функції створення ключа
\end{itemize}

\vspace{0.5em}


\end{frame}


\begin{frame}{Висновки}

\vspace{0.5em}
\begin{itemize}
  \item CtKDF \textbf{не є стійкою за визначенням}, якщо використовується слабка або передбачувана функція створення ключа
  \item Навіть при ідеальних схемах \textit{структура секрету може витікати}
  \item Побудована ігрова модель $G_0/G_1$ дозволяє \textbf{формально змоделювати атаку}
  \item Ігрова перевага зловмисника \textit{може бути близькою до 1}, що свідчить про повну втрату IND-CPA
\end{itemize}


\end{frame}

\begin{frame}
  \centering
  \vspace{2em}
  {\LARGE \textbf{Дякую за увагу!}} \\[2em]
  {\textbf{Запитання?}} \\[2em]


\end{frame}


\end{document}